\section{Elimination of normalized pair effects}
\label{sec:peps_pair_effects}
This section treats \cite[Lemma~5.1]{OpenAI2026PolynomialPEPS},
September 24, 2026, \texttt{04-compression.tex}, lines 53--127, together with
the choice of the stack length and the rescaling in the proof of
\cite[Theorem~5.2]{OpenAI2026PolynomialPEPS}, lines 201--212. The estimates
are proved for monomials written as chains of recorded pair effects between
arbitrary contractions, with no assignment of registers to parties; the
statements of the lemma about parties are discussed in
\cite{gap:polypeps_pair_effects_party_layout}.
% Maintainer note: the clauses on source-only expansions and on party
% ownership of the additional registers are not formalized; see the
% paper-gap note polypeps_pair_effects_party_layout.tex.

\begin{lemma}[Elimination of normalized pair effects]
    \label{lem:peps_pair_effects_source}
    \uses{thm:peps_pair_effects_operator,thm:peps_pair_combine}
    Let $G$ be a contraction with an expansion $G=\sum_\xi c_\xi M_\xi$ into
    $K$ allowed monomials on a fixed set of parties, each with at most $r$
    pair effects. For every integer $m\ge1$ there are a gate $\widetilde G_m$
    with additional garbage registers and a fixed normalized garbage vector
    $\Gamma_m$ with
    $\|\widetilde G_m-G\otimes|\Gamma_m\rangle\|\le(r/\sqrt m)\sum_\xi|c_\xi|$.
    The gate $\widetilde G_m$ has an expansion using only local contractions
    and normalized pair sources, with at most $Km^r$ monomials and absolute
    coefficient sum at most $\sum_\xi|c_\xi|$. Its additional registers are
    owned by the original participating parties, and all sources on the same
    pair of parties may be combined into one normalized pair source. This is
    \cite[Lemma~5.1]{OpenAI2026PolynomialPEPS}.
\end{lemma}

\begin{proof}
    \uses{thm:peps_pair_effects_operator,thm:peps_pair_combine}
    The error bound, the count and the coefficient sum are
    Theorem~\ref{thm:peps_pair_effects_operator}, and the combination of
    sources rests on Theorem~\ref{thm:peps_pair_combine}. The remaining clauses
    require placing the monomials on the parties, so that each piece between
    two effects is a product of local contractions and the stacks of an
    effect are held by its two endpoint parties.
\end{proof}

Throughout, $\eta$ is a unit vector of a finite-dimensional Hilbert space
$\mathcal W$, typically a pair space $\mathcal W=\mathcal U\otimes\mathcal V$,
and $m\ge1$.

\begin{definition}[Cyclic insertion map]
    \label{def:peps_cyclic_insertion}
    \lean{CyclicInsertion.insertAt}
    \lean{CyclicInsertion.cyclicInsertion}
    \lean{CyclicInsertion.idealInsertion}
    \leanok
    For a position $k\in\{1,\dots,m\}$, let $I_k:\mathcal W\to\mathcal W^{\otimes m}$
    place its input in register $k$ and a copy of $\eta$ in every other
    register. The cyclic insertion map is the average
    \[
        T_m=\frac1m\sum_{k=1}^m I_k,
    \]
    and the ideal insertion map is $P_m=|\eta^{\otimes m}\rangle\langle\eta|$.
\end{definition}

\begin{theorem}[Insertions as register permutations]
    \label{thm:peps_insertion_permutation}
    \lean{CyclicInsertion.insertAt_perm}
    \lean{CyclicInsertion.cyclicInsertion_eq_sum_addRight}
    \lean{CyclicInsertion.registerPerm_pair_index}
    \leanok
    \uses{def:peps_cyclic_insertion}
    For a permutation $\sigma$ of the registers, $I_{\sigma(k)}$ is $I_k$
    followed by the corresponding permutation of the tensor factors of
    $\mathcal W^{\otimes m}$. Hence $T_m$ is the average of the $m$ cyclic
    register permutations applied after the insertion in the first register.
    For a pair space $\mathcal W=\mathcal U\otimes\mathcal V$, the index map
    of a register permutation permutes the $\mathcal U$ indices and the
    $\mathcal V$ indices by the same permutation.
\end{theorem}

\begin{proof}\leanok
    Both sides agree on product vectors, coordinate by coordinate. The pair
    statement is the corresponding identity of index maps.
\end{proof}

\begin{theorem}[Contractivity and the fixed vector]
    \label{thm:peps_cyclic_insertion_basic}
    \lean{CyclicInsertion.norm_cyclicInsertion_le_one}
    \lean{CyclicInsertion.cyclicInsertion_self}
    \leanok
    \uses{def:peps_cyclic_insertion}
    The map $T_m$ is a contraction, and $T_m\eta=\eta^{\otimes m}$.
\end{theorem}

\begin{proof}\leanok
    Each $I_k$ is an isometry because $\|\eta\|=1$, so the average of the $m$
    maps has norm at most one. Each $I_k$ sends $\eta$ to
    $\eta^{\otimes m}$, and the average of $m$ equal terms is that term.
\end{proof}

\begin{theorem}[One-effect estimate]
    \label{thm:peps_one_effect_estimate}
    \lean{CyclicInsertion.norm_cyclicInsertion_of_inner_eq_zero}
    \lean{CyclicInsertion.norm_cyclicInsertion_sub_ideal_le}
    \lean{CyclicInsertion.norm_cyclicInsertion_sub_ideal_eq}
    \lean{CyclicInsertion.cyclicInsertion_sub_ideal_eq_zero}
    \lean{CyclicInsertion.norm_rTensor_cyclicInsertion_sub_ideal_le}
    \leanok
    \uses{def:peps_cyclic_insertion}
    For every $v\perp\eta$ one has $\|T_mv\|=\|v\|/\sqrt m$. Consequently
    $\|T_m-P_m\|\le m^{-1/2}$, with equality when the orthogonal complement of
    $\eta$ is nonzero, and $T_m=P_m$ when it is zero. The same upper bound
    holds for $(T_m-P_m)\otimes\mathbf 1_{\mathcal S}$ on any spectator space
    $\mathcal S$.
\end{theorem}

\begin{proof}\leanok
    \uses{thm:peps_cyclic_insertion_basic}
    For $v\perp\eta$, the vectors $I_jv$ and $I_kv$ with $j\ne k$ are
    orthogonal, since their inner product contains the factor
    $\langle v,\eta\rangle=0$, and each has norm $\|v\|$. Hence
    $\|\sum_kI_kv\|^2=m\|v\|^2$ and $\|T_mv\|=\|v\|/\sqrt m$. For a general
    input $v$, the difference $(T_m-P_m)v$ equals $T_m$ applied to the
    component $v-\langle\eta,v\rangle\eta$ orthogonal to $\eta$, whose norm is
    at most $\|v\|$. This gives the upper bound. A nonzero vector orthogonal to
    $\eta$ attains it, and when no such vector exists the difference vanishes.
    Tensoring with the identity does not increase the operator norm.
\end{proof}

\begin{definition}[Monomials with recorded pair effects]
    \label{def:peps_effect_chain}
    \lean{TNLean.PEPS.PairEffect.effectMap}
    \lean{TNLean.PEPS.PairEffect.EffectChain}
    \lean{TNLean.PEPS.PairEffect.EffectChain.eval}
    \lean{TNLean.PEPS.PairEffect.EffectChain.effectCount}
    \lean{TNLean.PEPS.PairEffect.EffectChain.IsAllowed}
    \leanok
    A monomial with $n$ recorded pair effects is a composition
    \[
        M=F_n\circ E_n\circ F_{n-1}\circ\cdots\circ E_1\circ F_0,
    \]
    where each $F_i$ ends with a pair register $\mathcal W_i$ beside a
    spectator space $\mathcal S_i$, and
    $E_i=\langle\eta_i|\otimes\mathbf 1_{\mathcal S_i}$ contracts that
    register with a pair bra. The monomial is allowed when every $F_i$ is a
    contraction and every $\eta_i$ is a unit vector. Here the pieces $F_i$
    are arbitrary contractions; no division of registers among parties is
    recorded.
\end{definition}

\begin{definition}[Replaced monomial and replaced gate]
    \label{def:peps_replaced_gate}
    \lean{TNLean.PEPS.PairEffect.EffectChain.stackVector}
    \lean{TNLean.PEPS.PairEffect.EffectChain.replace}
    \lean{TNLean.PEPS.PairEffect.EffectChain.replaceTerm}
    \lean{TNLean.PEPS.PairEffect.GateExpansion}
    \lean{TNLean.PEPS.PairEffect.gate}
    \lean{TNLean.PEPS.PairEffect.inventory}
    \lean{TNLean.PEPS.PairEffect.inventoryVector}
    \lean{TNLean.PEPS.PairEffect.replaceGate}
    \lean{TNLean.PEPS.PairEffect.termList}
    \leanok
    \uses{def:peps_cyclic_insertion,def:peps_effect_chain}
    Give each effect occurrence $E_i$ of a monomial $M$ its own output stack
    $\mathcal W_i^{\otimes m}$. The replaced monomial $\widetilde M_m$ replaces
    each $E_i$ by $T_m$ for $\eta_i$, writing into the stack of that
    occurrence, which stays idle afterwards. Its ideal stack content is
    $\bigotimes_i\eta_i^{\otimes m}$. For a gate expansion
    $G=\sum_\xi c_\xi M_\xi$, the stack space of $G$ is the tensor product of
    the stacks of all effect occurrences of all monomials, and
    $\Gamma_m=\bigotimes_o\eta_o^{\otimes m}$ is the product of all ideal stack
    contents. The replaced gate $\widetilde G_m$ is $\sum_\xi c_\xi$ times the
    branch in which $\widetilde M_{\xi,m}$ acts and the stacks of every other
    monomial are prepared in their ideal content. Selecting one insertion
    position $\kappa_i$ for each occurrence in place of the average defines one
    term of the expansion of $\widetilde M_m$, with weight $m^{-n}$.
\end{definition}

\begin{theorem}[Error of one replaced monomial]
    \label{thm:peps_replaced_monomial_error}
    \lean{TNLean.PEPS.PairEffect.EffectChain.norm_replace_sub_le}
    \lean{TNLean.PEPS.PairEffect.EffectChain.norm_eval_le_one}
    \lean{TNLean.PEPS.PairEffect.EffectChain.norm_stackVector}
    \lean{TNLean.PEPS.PairEffect.EffectChain.norm_replace_le_one}
    \leanok
    \uses{def:peps_replaced_gate}
    For an allowed monomial $M$ with $n$ recorded pair effects, $M$ and
    $\widetilde M_m$ are contractions, the ideal stack content is a unit
    vector, and
    \[
        \Bigl\|\widetilde M_m-M\otimes\Bigl|\bigotimes_i\eta_i^{\otimes m}\Bigr\rangle\Bigr\|
        \le\frac n{\sqrt m}.
    \]
\end{theorem}

\begin{proof}\leanok
    \uses{thm:peps_cyclic_insertion_basic,thm:peps_one_effect_estimate}
    By induction on $n$. Appending the ideal content of the first stack after
    $E_1$ is the same as applying $P_m\otimes\mathbf 1$ in place of $E_1$. The
    difference therefore splits into the replacement error of the remaining
    monomial composed with $T_m\otimes\mathbf 1$, and the ideal remainder
    composed with $(T_m-P_m)\otimes\mathbf 1$. All other factors are
    contractions, so the two terms contribute $(n-1)/\sqrt m$ and $1/\sqrt m$.
\end{proof}

\begin{theorem}[Expansion of a replaced monomial]
    \label{thm:peps_replaced_monomial_expansion}
    \lean{TNLean.PEPS.PairEffect.EffectChain.replace_eq_sum}
    \lean{TNLean.PEPS.PairEffect.EffectChain.norm_replaceTerm_le_one}
    \leanok
    \uses{def:peps_replaced_gate}
    The replaced monomial with $n$ recorded pair effects equals
    $m^{-n}\sum_\kappa\widetilde M_{m,\kappa}$, the sum over all choices
    $\kappa$ of one insertion position per occurrence. Every term
    $\widetilde M_{m,\kappa}$ is a contraction.
\end{theorem}

\begin{proof}\leanok
    \uses{def:peps_cyclic_insertion}
    Expand each average $T_m$ by linearity and collect the $m^n$ choices of
    positions, by induction on $n$. Each term composes the contractions
    $F_i$ with the isometries $I_{\kappa_i}$ and tensor rearrangements.
\end{proof}

\begin{theorem}[Pair-effect elimination, operator estimates]
    \label{thm:peps_pair_effects_operator}
    \lean{TNLean.PEPS.PairEffect.pairEffectElimination}
    \lean{TNLean.PEPS.PairEffect.norm_inventoryVector}
    \lean{TNLean.PEPS.PairEffect.norm_replaceGate_sub_le}
    \lean{TNLean.PEPS.PairEffect.replaceGate_eq_sum_termList}
    \lean{TNLean.PEPS.PairEffect.length_termList_le}
    \lean{TNLean.PEPS.PairEffect.sum_norm_coeff_termList}
    \lean{TNLean.PEPS.PairEffect.norm_term_le_one}
    \leanok
    \uses{def:peps_replaced_gate}
    Let $G=\sum_\xi c_\xi M_\xi$ be an expansion into $K$ allowed monomials,
    each with at most $r$ recorded pair effects, and let $m\ge1$. Then
    $\Gamma_m$ is a unit vector,
    \[
        \bigl\|\widetilde G_m-G\otimes|\Gamma_m\rangle\bigr\|
        \le\frac r{\sqrt m}\sum_\xi|c_\xi|,
    \]
    and $\widetilde G_m$ is a weighted sum of at most $Km^r$ contractions
    whose absolute coefficient sum equals $\sum_\xi|c_\xi|$.

    These are the estimates of \cite[Lemma~5.1]{OpenAI2026PolynomialPEPS}
    for monomials given as chains of recorded pair effects: the error bound,
    the common garbage vector, the monomial count and the coefficient sum.
    The statements about local contractions, normalized pair sources and
    ownership of the additional registers by the parties are those of
    Lemma~\ref{lem:peps_pair_effects_source}, see
    \cite{gap:polypeps_pair_effects_party_layout}. The hypothesis that $G$ is
    a contraction is not needed.
    % Maintainer note: only the operator estimates are formalized here.
\end{theorem}

\begin{proof}\leanok
    \uses{thm:peps_replaced_monomial_error,thm:peps_replaced_monomial_expansion}
    By induction on the list of monomials. The difference splits into $c_\xi$
    times the error of the first replaced monomial, followed by appending the
    unit vector of the remaining stacks, plus the error of the remaining gate
    with the first stacks appended. The first term is bounded by
    Theorem~\ref{thm:peps_replaced_monomial_error} and the second by induction;
    appending a unit vector and rearranging tensor factors are contractions.
    Expanding each replaced monomial by
    Theorem~\ref{thm:peps_replaced_monomial_expansion} gives $m^{n_\xi}\le m^r$
    terms of coefficient $c_\xi m^{-n_\xi}$ for the monomial $\xi$, whose
    absolute values sum to $|c_\xi|$.
\end{proof}

\begin{theorem}[Combining pair sources]
    \label{thm:peps_pair_combine}
    \lean{EuclideanSpace.pairCombine}
    \lean{EuclideanSpace.norm_pairCombine}
    \leanok
    For pair vectors $\eta\in\mathcal U\otimes\mathcal V$ and
    $\eta'\in\mathcal U'\otimes\mathcal V'$, regard $\eta\otimes\eta'$ as a
    pair vector in $(\mathcal U\otimes\mathcal U')\otimes(\mathcal V\otimes\mathcal V')$.
    Its norm is $\|\eta\|\,\|\eta'\|$; in particular, normalized pair vectors
    combine into a normalized pair vector.
\end{theorem}

\begin{proof}\leanok
    The squared norm is a sum of products of squared coordinates, which
    factorizes after regrouping the index set.
\end{proof}

\begin{corollary}[Stack length and rescaling]
    \label{cor:peps_pair_effects_rescaling}
    \lean{TNLean.PEPS.PairEffect.stackLength}
    \lean{TNLean.PEPS.PairEffect.stackLength_spec}
    \lean{TNLean.PEPS.PairEffect.replaceGate_rescaled}
    \lean{TNLean.PEPS.PairEffect.replaceGate_rescaled_stackLength}
    \leanok
    \uses{thm:peps_pair_effects_operator}
    In the setting of Theorem~\ref{thm:peps_pair_effects_operator}, suppose
    moreover that $G$ is a contraction, and let $\delta>0$ and
    $S\ge\sum_\xi|c_\xi|$. With $m=\lceil(rS/\delta)^2\rceil+1$, the replaced
    gate satisfies $\|\widetilde G_m-G\otimes|\Gamma_m\rangle\|\le\delta$ and
    $\|\widetilde G_m\|\le1+\delta$, and $\widetilde G_m/(1+\delta)$ is a
    contraction within $2\delta$ of $G\otimes|\Gamma_m\rangle$. The number $m$
    depends only on $r$, $S$ and $\delta$.
\end{corollary}

\begin{proof}\leanok
    The choice of $m$ gives $\sqrt m\ge rS/\delta$, so the bound of
    Theorem~\ref{thm:peps_pair_effects_operator} is at most $\delta$. Since
    $G\otimes|\Gamma_m\rangle$ is a contraction, the triangle inequality gives
    $\|\widetilde G_m\|\le1+\delta$. Dividing by $1+\delta$ yields a
    contraction, and its distance to $G\otimes|\Gamma_m\rangle$ is at most
    $\delta/(1+\delta)+\delta\le2\delta$.
\end{proof}
